\documentclass{article}
\usepackage[utf8]{inputenc}
\usepackage{a4wide}

\begin{document}

\section*{นักตกปลากับแหล่งน้ำปริศนา}

ธีรภพได้รับแผนที่พื้นที่แห่งหนึ่งในรูปแบบของ ตารางขนาดm x n

(matrix)โดยในแต่ละช่องของตารางจะมีค่าเป็นเลขจำนวนเต็มgrid{[}r{]}{[}c{]}ซึ่งหมายถึง:\\



\begin{itemize}

\tightlist

\item

  ถ้าgrid{[}r{]}{[}c{]} == 0ช่องนั้นคือ พื้นดิน ไม่สามารถตกปลาได้

\item

  ถ้าgrid{[}r{]}{[}c{]} \textgreater{} 0ช่องนั้นคือ แหล่งน้ำ

  ที่มีปลาจำนวนgrid{[}r{]}{[}c{]}ตัว

\end{itemize}



\section{\texorpdfstring{\textbf{กฎของเกม:}}{กฎของเกม:}}\label{uxe01uxe0euxe02uxe2duxe07uxe40uxe01uxe21}



ธีรภพสามารถเริ่มต้นที่ จุดใดก็ได้ ที่เป็นแหล่งน้ำ แล้วทำสิ่งต่อไปนี้ได้ไม่จำกัดจำนวนครั้ง:



1.ตกปลาทั้งหมด ที่อยู่ในช่องปัจจุบัน (ปลาหมดแล้วจะเป็นแหล่งน้ำว่างเปล่า)



2.ย้ายไปช่องใกล้เคียง ที่เป็นแหล่งน้ำ (ห้ามเดินลงดิน)



โดย ช่องใกล้เคียง (adjacent)คือช่องที่อยู่:



\begin{itemize}

\tightlist

\item

  ขวา(r, c + 1)

\item

  ซ้าย(r, c - 1)

\item

  ล่าง(r + 1, c)

\item

  บน(r - 1, c)

\end{itemize}



\subsection{\texorpdfstring{\textbf{เป้าหมายของคุณ:}}{เป้าหมายของคุณ:}}\label{uxe40uxe1buxe32uxe2buxe21uxe32uxe22uxe02uxe2duxe07uxe04uxe13}



จงหาว่า ธีรภพสามารถตกปลาได้มากที่สุดกี่ตัว หากเขาเลือกจุดเริ่มต้นได้อย่างดีที่สุด\\

หากไม่มีแหล่งน้ำเลย ให้ตอบ0



\subsection{\texorpdfstring{\textbf{ข้อจำกัด:}}{ข้อจำกัด:}}\label{uxe02uxe2duxe08uxe33uxe01uxe14}



\begin{itemize}

\tightlist

\item

  ขนาดของแผนที่:1 \textless= m, n \textless= 10

\item

  จำนวนปลาสูงสุดในแต่ละช่อง:0 \textless= grid{[}i{]}{[}j{]} \textless= 10

\end{itemize}



\paragraph{\texorpdfstring{\textbf{ตัวอย่างที่1:}}{ตัวอย่างที่1:}}\label{uxe15uxe27uxe2duxe22uxe32uxe07uxe171}



\textbf{Input:}



\textbf{\includegraphics[width=1.55208in,height=1.0625in,alt={image.png}]{/public/upload/ca5b7d4622.png}\\

}



grid = {[}{[}0,2,1,0{]},



{[}4,0,0,3{]},



{[}1,0,0,4{]},



{[}0,3,2,0{]}{]}



\textbf{Output:}7\\

อธิบาย: เริ่มที่(1,3)ได้3ตัว เดินลงมาที่(2,3)ได้อีก4ตัว รวม7ตัว



\hfill\break



\paragraph{\texorpdfstring{\textbf{ตัวอย่างที่2:}}{ตัวอย่างที่2:}}\label{uxe15uxe27uxe2duxe22uxe32uxe07uxe172}



\textbf{Input:}



grid = {[}{[}1,0,0,0{]},



{[}0,0,0,0{]},



{[}0,0,0,0{]},



{[}0,0,0,1{]}{]}



\textbf{Output:}\texttt{1}\\

อธิบาย: จุดที่มีปลาอยู่คือ\texttt{(0,0)}และ\texttt{(3,3)}ซึ่งต่างคนต่างอยู่

ไม่สามารถเดินไปที่อื่นได้



\subsection*{Input}
บรรทัดที่1 : ขนาดของ m และ n คั่นด้วยช่องว่าง



บรรทัดที่ 2 - m+1 : ตัวเลขในช่องตารางจำนวน n คั่นด้วยช่องว่าง



\subsection*{Output}
บรรทัดที่ 1 : จำนวนปลาทีตกได้สูงสุด~หากเขาเลือกจุดเริ่มต้นได้อย่างดีที่สุด

(หากไม่มีแหล่งน้ำเลย ให้ตอบ0 )



\end{document}
